This element simulates a wiggler or undulator using a modified version of Ying Wu's
canonical integration code for wigglers.  To use the element, one must
supply an SDDS file giving harmonic analysis of the wiggler field.
The field expansion used by the code for a horizontally-deflecting
wiggler is (Y. Wu, Duke University, private communication).
\begin{equation}
B_y = -\left|B_0\right| \sum_{m,n} C_{mn}\cos(k_{xl} x) \cosh (k_{ym} y)
\cos(k_{zn} z + \theta_{zn}),
\end{equation}
where $\left|B_0\right|$ is the peak value of the on-axis magnetic field,
the $C_{mn}$ give the relative amplitudes of the harmonics, the wavenumbers
statisfy $k^2_{ym} = k^2_{xl} + k^2_{zn}$, and $\theta_{zn}$ is the phase.

The file must contain the following columns:
\begin{itemize}
\item The harmonic amplitude, $C_{mn}$, in column {\tt Cmn}.
\item The phase, in radians, in column {\tt Phase}.  The phase of the first
harmonic should be 0 or $\pi$ in order to have matched dispersion.
\item The three wave numbers, normalized to $k_w = 2\pi/\lambda_w$, where
 $\lambda_w$ is the wiggler period.  These are given in columns 
 {\tt KxOverKw}, {\tt KyOverKw}, and {\tt KzOverKw}.
\end{itemize}

For matrix-based computations {\tt elegant} uses a first-order matrix
derived from particle tracking when it encounterse a CWIGGLER.  Tests
show that this gives good agreement in the tunes from tracking and
Twiss parameter calculations.  For radiation integrals, an idealized
sinusoidal wiggler model is used with bending radius equal to
$B\rho/(B_0 \sum C_{mn})$ for each plane.  Energy loss, energy spread,
and horizontal emittance should be estimated accurately.

{\tt elegant} allows specifying field expansions for on-axis $B_y$ and
$B_x$ components, so one can model a helical wiggler.  However, in this
case one set of components should have $\theta_{zn} = 0$ or $\theta_{zn}=\pi$,
while the other should have $\theta_{zn} = \pm \pi/2$.
Using Wu's code, the latter set will not have matched dispersion.  Our
modified version solves this by delaying the beginning of the field
components in question by $\lambda/4$ and ending the field prematurely
by $3\lambda/4$.  This causes all the fields to start and end at the
crest, which ensures matched dispersion.  The downside is that the
(typically) vertical wiggler component is missing a full period of
field.  One can turn off this behavior by setting \verb|FORCE_MATCHED=0|.

{\bf Additional field expansions}

Y. Wu's code included field expansions for a vertically-deflecting
wiggler as well as the horizontally-deflecting wiggler given above.
In both cases, these expansions are suitable for a wiggler with two
poles that are above/below or left/right of the beam axis.  They are
not always suitable for devices with more complex pole geometries.

Another geometry that is important is a ``split pole'' wiggler, in which
each pole is made from two pieces.  Such configurations are seen, for example,
in devices used to produce variable polarization.  In such cases, the
expansion given above may not be appropriate.  Here, we summarize the form of the
various expansions that {\tt elegant} supports.  For brevity, we show the
form of a single harmonic component.

Horizontal wiggler, normal poles, produces $B_y$ only on-axis.   Specified
by setting {\tt BY\_SPLIT\_POLE=0}, and giving {\tt BY\_FILE} or {\tt SINUSOIDAL=1} with {\tt VERTICAL=0}.
\begin{eqnarray}
B_x & = & \left|B_0\right| \frac{k_x \cos (k_z z + \phi) \sin (k_x x) \sinh (k_y y)}{k_y}\\
B_y & = & - \left|B_0\right| \cos (k_x x) \cos (k_z z + \phi) \cosh (k_y y)
\end{eqnarray}
where $k_y^2 = k_x^2 + k_z^2$.

{\bf Experimental feature:} Horizontal wiggler, normal poles, with transverse gradient, produces $B_y$ only on-axis.   Specified
by setting {\tt BY\_SPLIT\_POLE=0}, {\tt SINUSOIDAL=1}, {\tt TGU=1}, {\tt VERTICAL=0}.
The TGU normalized gradient is given using the \verb|TGU_GRADIENT| parameter.
Taking $a$ as the normalized gradient, the fields are\cite{RLindbergPC}
\begin{eqnarray}
B_x & = & \frac{a \left|B_0\right| \sinh k_u y \cos k_u z}{k_u} \\
B_y & = & \left|B_0\right| \left((1 + a x) \cosh k_u y \cos k_u z + \frac{a C }{2 k_u^2}\frac{e \left|B_0\right|}{\gamma m_e c }\right) \\
B_z & = & -\left|B_0\right| (1 + a x) \sinh k_u y \sin k_u z,
\end{eqnarray}
where $k_u = k_y = k_z$, $k_x = 0$ is assumed, $\gamma$ is the central relativistic factor for the beam and $C$ is given by the
\verb|TGU_COMP_FACTOR| parameter.
This factor, and the term it multiplies, is present in order to help suppress the 
trajectory error at the end of the device.
It may require adjustment in order to achieve the desired level of correction.
In addition, the user may need to adjust the pole-strength factors and include external misalignments
and steering magnets in order to supress not only the trajectory error, but also dispersion errors.

Horizontal wiggler, split poles, produces $B_y$ only on-axis.   Specified
by setting {\tt BY\_SPLIT\_POLE=1}, and giving {\tt BY\_FILE} or {\tt SINUSOIDAL=1} with {\tt VERTICAL=0}.
\begin{eqnarray}
B_x & = & -\left|B_0\right| \frac{k_x \cos (k_z z + \phi) \sin (k_y y) \sinh (k_x x)}{k_y}\\
B_y & = & -\left|B_0\right| \cos (k_y y) \cos (k_z z + \phi) \cosh (k_x x)
\end{eqnarray}
where $k_x^2 = k_y^2 + k_z^2$.

Vertical wiggler, normal poles, produces $B_x$ only on-axis.   Specified
by setting {\tt BX\_SPLIT\_POLE=0}, and giving {\tt BX\_FILE} or {\tt SINUSOIDAL=1} with either {\tt VERTICAL=1} or
{\tt HELICAL=1}.
\begin{eqnarray}
B_x & = & \left|B_0\right| \cos (k_y y) \cos (k_z z + \phi) \cosh (k_x x)\\
B_y & = & -\left|B_0\right| \frac{k_y \cos (k_z z + \phi) \sin (k_y y) \sinh (k_x x)}{k_x}
\end{eqnarray}
where $k_x^2 = k_y^2 + k_z^2$.

Vertical wiggler, split poles, produces $B_x$ only on-axis.   Specified
by setting {\tt BX\_SPLIT\_POLE=1}, and giving {\tt BX\_FILE} or {\tt SINUSOIDAL=1} with either {\tt VERTICAL=1} or
{\tt HELICAL=1}.
\begin{eqnarray}
B_x & = & \left|B_0\right| \cos (k_x x) \cos (k_z z + \phi) \cosh (k_y y)\\
B_y & = & \left|B_0\right| \frac{k_y \cos (k_z z + \phi) \sin (k_x x) \sinh (k_y y)}{k_x}
\end{eqnarray}
where $k_y^2 = k_x^2 + k_z^2$.

{\bf Splitting wigglers}

The \verb|CWIGGLER| element supports a limited ability to split a long element into parts using the
\verb|element_divisions| command or the \verb|divide_elements| pararameter of the \verb|run_setup| command.

In addition, if contiguous \verb|CWIGGLER| elements are seen in a beamline, they will be treated as part of the
same element. That means that the pole factors will be ignored except at the ends of the sequence.
For this purpose, \verb|CWIGGLER| elements separated only by \verb|MARK| or \verb|WATCH| elements are considered
to be contiguous. 

There are several ways to specify the strength of the magnetic field:
\begin{enumerate}
\item Using the \verb|BMAX|, \verb|BYMAX|, or \verb|BXMAX| parameters directly.
\item Using a predefined undulator model to compute the value of \verb|BYMAX|. Several hybrid-permanent-magnet undulator (HPMU) models are supported. To use these, the \verb|GAP| parameter
      must be used to provide the full magnetic gap.
      \begin{itemize}
        \item \verb|Halbach| --- The original model of a hybrid-permanent-magnet (HPM) undulator by K. Halbach \cite{HalbachModel}. This is mostly of historical interest.
        \item \verb|SmCo-RT| --- A SmCo HPMU at room temperature, using the model of R. Dejus \cite{DejusModels}.
        \item \verb|NdFeB-RT| --- A NdFeB HPMU at room temperature, using the model of R. Dejus \cite{DejusModels}.
        \item \verb|NdFeB-150K| --- A NdFeB HPMU at 150 K, using the model of E. Moog \cite{MoogModels}.
        \item \verb|PrFeB-77K| --- A PrFeB HPMU at 77 K, using the model of E. Moog \cite{MoogModels}.
        \item \verb|Custom| --- A user-specified Halbach-type model. The parameters \verb|C1|, \verb|C2|, and \verb|C3| must be given. The field is computed from
        \begin{equation}
                B = C_1 e^{C_2 r + C_3 r^3},
        \end{equation}
        where $r = g/\lambda_u$, $g$ is the full magnetic gap, and $\lambda$ is the undulator period.
      \end{itemize}
      Two superconducting undulator (SCU) models are supported. In addition to providing the magnetic gap, the user must
      provide a value for the parameter \verb|JFRACTION|, which is the fraction of the quench current at which to
      operate; 0.8 is considered conservative. These are based on the model of S. H. Kim \cite{Kim-SCU-Scaling}.
      \begin{itemize}
        \item \verb|NbTi| --- A NbTi SCU.
        \item \verb|Nb3Sn| --- A ${\rm Nb_3 Sn}$ SCU.  This is just the \verb|NbTi| model scaled up by 30\% \cite{IvanyushenkovPC}
      \end{itemize}
\end{enumerate}
